\section{Importance of Software Synthesis at Processor Level}
Software synthesis produces binary codes for programmable processing elements~\cite{Gajski.2009}. As shown in Figure~\ref{fig:Software_Synthesis}, a system model of a processor describes the components and the connectivity, whereas the bottom portion of the design flow takes the system model as an input and generates an implementation of each component~\cite{Gajski.2009}.
\begin{figure}[htbp]
    \centering
    \includegraphics[width=0.9\columnwidth]{C:/Users/exone/Documents/GitHub/Hardware_Software_Co-design_seminar/Sections/Diagrams/Software_Synthesis.png}
    \caption{Software Synthesis~\cite{Gajski.2009}.}
    \label{fig:Software_Synthesis}
\end{figure}  
Software development for embedded system covers most of the expeses of producing modern multi-processor systems in compared to hardware development. As there are more capabilities of any processor these days, they need firmware, drivers, OS code among others to work properly. Hence, lies the development of multi-faceted software programs to embed in one processor. However, hand-written codes, like conventional practice, do take lot of time for implementation. Hence, engineers, these days, prefer to implement features of processors in software tools because of their flexibility, reusability and for being easily updatable.

